\begin{theorem}[SmoothedChebyshevClose]
\uses{SmoothedChebyshev}
\label{SmoothedChebyshevClose}
\lean{SmoothedChebyshevClose}
% at /home/runner/work/trough/trough/library/PrimeNumberTheoremAnd/MediumPNT.lean:689.0-805.23
\leanok
We have that
  $$\psi_{\epsilon}(X) = \psi(X) + O(\epsilon X \log X).$$
\end{theorem}
\begin{proof}
\uses{Smooth1Nonneg,SmoothedChebyshevDirichlet,Smooth1Properties_above,Smooth1Properties_below,Smooth1,Smooth1LeOne}
\leanok
Take the difference. By Lemma \ref{Smooth1Properties_above} and \ref{Smooth1Properties_below},
  the sums agree except when $1-c \epsilon \leq n/X \leq 1+c \epsilon$. This is an interval of
  length $\ll \epsilon X$, and the summands are bounded by $\Lambda(n) \ll \log X$.

  %[No longer relevant, as we will do better than any power of log savings...: This is not enough,
  %as it loses a log! (Which is fine if our target is the strong PNT, with
  %exp-root-log savings, but not here with the ``softer'' approach.) So we will need something like
  %the Selberg sieve (already in Mathlib? Or close?) to conclude that the number of primes in this
  %interval is $\ll \epsilon X / \log X + 1$.
  %(The number of prime powers is $\ll X^{1/2}$.)
  %And multiplying that by $\Lambda (n) \ll \log X$ gives the desired bound.]
\end{proof}
